\documentclass[11pt]{article}
\usepackage{amsmath}
\usepackage{amssymb}
\usepackage{amsthm}
\usepackage[margin=1in]{geometry}

\newtheorem{theorem}{Theorem}
\newtheorem{proposition}[theorem]{Proposition}
\newtheorem{lemma}[theorem]{Lemma}
\newtheorem{corollary}[theorem]{Corollary}
\theoremstyle{definition}
\newtheorem{definition}[theorem]{Definition}
\newtheorem{remark}[theorem]{Remark}

\newcommand{\Zp}{\mathbb{Z}_p}
\newcommand{\D}{\Delta}
\newcommand{\vpp}{\nu_p}

\title{A disproof of conjecture 00000000750:\\
the exact eventual period of the reduced Mahler coefficient\\
sequence is not $p^{k-1}(p-1)$}
\author{earthking11}
\date{13 September 2026}

\begin{document}
\maketitle

\begin{abstract}
We refute conjecture 00000000750. The conjecture asserts that for the
$\bmod\,p^k$ reduction of a $1$-Lipschitz map on $\Zp$, the Mahler coefficient
sequence $(a_n \bmod p^k)$ is eventually periodic with \emph{exact} period
$p^{k-1}(p-1)$, ``which cannot be shortened''. We exhibit the map
$f(x)=x+1$, which is an isometry of $\Zp$ (hence $1$-Lipschitz), is
Haar-measure-preserving, and is minimal. Its Mahler coefficient sequence is
\[
(a_n)_{n\ge 0} = (1,1,0,0,0,\dots).
\]
Reduced modulo $p^k$ this sequence is eventually constant $0$, so its minimal
eventual period is $1$. For the test pairs $(p,k) \in \{(2,2),(2,3),(3,1),
(3,2),(5,1),(5,2)\}$ the conjectured value $p^{k-1}(p-1)$ is respectively
$2,4,2,6,4,20$, never $1$. Hence the conjecture is false. We also record that
the conjecture's quotation of Anashin's criterion appears to be misstated: the
correct divisibility condition is $p^{\vpp(n!)} \mid a_n$, not
$p^n \mid a_n$; the quoted form already fails for $f(x)=x+1$ because
$p \nmid a_1 = 1$ whenever $p>1$.
\end{abstract}

\section{Setup: Mahler coefficients}

Let $p$ be a prime and let $f:\Zp \to \Zp$ be continuous. Every such function
has a unique Mahler expansion
\begin{equation}\label{eq:mahler}
  f(x) = \sum_{n\ge 0} a_n \binom{x}{n},
  \qquad
  a_n = (\D^n f)(0) = \sum_{i=0}^{n} (-1)^{n-i}\binom{n}{i} f(i),
\end{equation}
where $\D$ is the forward difference operator
\[
  (\D f)(x) = f(x+1) - f(x),
  \qquad
  (\D^{n+1} f) = \D(\D^n f).
\]
The numbers $a_n \in \Zp$ are the \emph{Mahler coefficients} of $f$.

\begin{remark}[Anashin's criterion]\label{rem:anashin}
A standard form of Anashin's criterion states that a continuous
$f:\Zp\to\Zp$ is $1$-Lipschitz (equivalently, the coefficients satisfy the
compatibility needed for the expansion to converge) precisely when
\begin{equation}\label{eq:anashin}
  p^{\vpp(n!)} \mid a_n \quad \text{for all } n \ge 0,
\end{equation}
where $\vpp(n!) = \sum_{j\ge 1}\lfloor n/p^j\rfloor$. The conjecture under
consideration quotes this condition as ``$p^n \mid a_n$''. That quoted form is
not equivalent to \eqref{eq:anashin}: for $f(x)=x+1$ one has $a_1 = 1$, so
$p \mid a_1$ fails for every $p>1$, even though $f$ is an isometry and
certainly $1$-Lipschitz. This is a supporting observation about the wording of
the conjecture; the refutation below does not rely on it.
\end{remark}

\section{The counterexample $f(x)=x+1$}

Work with $f(x) = x+1$ on $\Zp$.

\begin{lemma}\label{lem:properties}
The map $f(x)=x+1$ has the following properties.
\begin{enumerate}
  \item It is an isometry of $\Zp$: $|f(x)-f(y)|_p = |x-y|_p$ for all
        $x,y \in \Zp$. In particular it is $1$-Lipschitz.
  \item It preserves Haar measure (translation is measure-preserving).
  \item It is minimal: the orbit $\{f^{\circ n}(0) : n \ge 0\} = \mathbb{N}
        = \{0,1,2,\dots\}$ is dense in $\Zp$.
\end{enumerate}
\end{lemma}

\begin{proof}
(1) $f(x)-f(y) = x-y$, so the $p$-adic absolute values agree. (2) $f$ is a
translation, and Haar measure on the locally compact group $\Zp$ is translation
invariant. (3) $f^{\circ n}(0) = n$, and the natural numbers are dense in
$\Zp$ (the residues $\bmod\,p^m$ exhaust every ball).
\end{proof}

\begin{proposition}[Mahler coefficients]\label{prop:coeffs}
For $f(x)=x+1$ the Mahler coefficient sequence is
\[
  (a_0,a_1,a_2,a_3,\dots) = (1,1,0,0,0,\dots).
\]
Equivalently, $(\D^n f)(0) = 0$ for every $n \ge 2$.
\end{proposition}

\begin{proof}
First, $\D(x+1) = 1$ is the constant function $1$:
\[
  (\D f)(x) = f(x+1)-f(x) = (x+2)-(x+1) = 1 .
\]
Hence $\D^n(x+1) = 0$ for every $n \ge 1$ applied to this constant, i.e.
$\D^2 f = \D 1 = 0$ and all higher differences vanish. Evaluating at $x=0$:
\[
  a_0 = f(0) = 1, \qquad a_1 = (\D f)(0) = 1, \qquad
  a_n = (\D^n f)(0) = 0 \ \ (n\ge 2).
\]
Equivalently, the Mahler expansion of $f$ terminates:
$x+1 = \binom{x}{0} + \binom{x}{1}$, matching $a_0=a_1=1$ and $a_n=0$ for
$n\ge 2$.
\end{proof}

\section{Periodicity modulo $p^k$}

Fix a prime $p$ and an integer $k \ge 1$, and reduce the sequence
$(a_n)$ modulo $p^k$.

\begin{theorem}\label{thm:period}
The reduced sequence $(a_n \bmod p^k)_{n\ge 0}$ is
\[
  (1,1,0,0,0,\dots) \bmod p^k ,
\]
which is eventually constant (equal to $0$). Consequently its minimal eventual
period is
\[
  T_{\mathrm{actual}} = 1 .
\]
\end{theorem}

\begin{proof}
Since $p^k \ge 2$, we have $1 \bmod p^k = 1 \ne 0$, so the first two terms
remain $1,1$. By Proposition~\ref{prop:coeffs}, $a_n = 0$ for all $n \ge 2$,
so $a_n \bmod p^k = 0$ for all $n\ge 2$; the tail is constant, and the smallest
positive eventual period of a constant tail is $1$.
\end{proof}

\begin{theorem}[Refutation of conjecture 00000000750]\label{thm:refute}
The conjectured exact period is
\[
  T_{\mathrm{claimed}} = p^{k-1}(p-1).
\]
For the test pairs
\[
  (p,k) \in \{(2,2),(2,3),(3,1),(3,2),(5,1),(5,2)\}
\]
the claimed value equals
\[
  2,\ 4,\ 2,\ 6,\ 4,\ 20
\]
respectively, whereas the actual minimal eventual period is $1$ in every case.
More generally, for $p>1$ and $k\ge 1$ one has $p^{k-1}(p-1) = 1$ only for
$(p,k)=(2,1)$, so for every pair with $(p,k)\ne(2,1)$ the claimed period
differs from the actual period $1$. Therefore the assertion that the exact
period is $p^{k-1}(p-1)$ and ``cannot be shortened'' fails.
\end{theorem}

\begin{proof}
By Theorem~\ref{thm:period} the actual period is $1$. The listed values of
$p^{k-1}(p-1)$ are obtained by direct computation:
$2^1\cdot1=2$, $2^2\cdot1=4$, $3^0\cdot2=2$, $3^1\cdot2=6$,
$5^0\cdot4=4$, $5^1\cdot4=20$. None equals $1$. For the general claim:
$p^{k-1}\ge 1$ and $p-1\ge 1$ for $p>1$, so $p^{k-1}(p-1)\ge p-1 \ge 1$; if
$p\ge 3$ then $p-1\ge 2$, and if $p=2$ then $p-1=1$ and $p^{k-1}(p-1)=2^{k-1}$,
which equals $1$ exactly when $k-1=0$, i.e.\ $k=1$. Thus the unique solution of
$p^{k-1}(p-1)=1$ with $p>1$, $k\ge1$ is $(2,1)$.
\end{proof}

\begin{remark}
The hypotheses of the conjecture are satisfied by the counterexample: by
Lemma~\ref{lem:properties}, $f(x)=x+1$ is $1$-Lipschitz, Haar-measure
preserving, and minimal. The conjecture asserts that the exact period is
attained ``simultaneously with the measure-preserving property of minimal
maps''; $f$ is such a map, yet its coefficient sequence has eventual period
$1$. Hence the universal claim is false.
\end{remark}

\section*{Verification summary}
\begin{itemize}
  \item \texttt{lean4/Main.lean} formalises the finite-difference core in core
        Lean (only \texttt{Int}/\texttt{Nat}, \texttt{import Std}, no Mathlib):
        the first difference of $x+1$ is the constant $1$, the second and all
        higher differences vanish, hence the Mahler coefficients are
        $(1,1,0,0,\dots)$; it also proves by \texttt{decide} that
        $p^{k-1}(p-1)\ne 1$ for the listed pairs.
  \item \texttt{reproduce.py} computes the Mahler coefficients by finite
        differences and compares the exact eventual period with the claimed
        value for the listed $(p,k)$; it prints \texttt{PASS}.
  \item \texttt{lean4/Check.lean} prints the axiom dependencies of each
        theorem.
\end{itemize}

\end{document}
